\section{Exactly consistent rounding for the full preordering}\label{sec:kernels}

This section proves that every feasible family of the sparse box preordering
of order $r$ can be rounded to a probability law on the box whose expected
objective exceeds the relaxation objective by at most $3\Abud/(2m^2+1)$,
where $m=\lfloor r/w\rfloor+1$. Consequently
$\fstar-\rhopre_r\le3w^2\Abud/(2r^2)$. The inverse-square rate itself is not
new (\cref{rem:pre-prior}). We give a complete proof because its construction
is the template for every upper bound in the paper, and because it yields
explicit constants, a finite-grid rounding, and a real certificate.

The construction applies one fixed univariate polynomial kernel $\Kj(x,u)$
in every coordinate of every bag and evaluates the local functional on the
resulting tensor kernel:
\[
  h_b(u)=L_b\Bigl(\prod_{i\in B_b}\Kj(x_i,u_i)\Bigr),
  \qquad u\in[-1,1]^{B_b}.
\]
Four properties of the kernel do four different jobs.
\begin{enumerate}[label=(\arabic*)]
\item \emph{Finite degree.} The kernel has degree $2(m-1)$ in $x_i$, so
  the tensor kernel lies in the domain of $L_b$.
\item \emph{Certified positivity.} For each fixed output $u_i$, the
  polynomial $\Kj(\cdot,u_i)$ is nonnegative on $[-1,1]$, and its interval
  certificate has the right degree. Multiplying the certificates over a bag
  gives an element of the truncated preordering, so $h_b\ge0$ although $L_b$
  need not be a measure.
\item \emph{Exact mass.} $\int\Kj(x,u)\,d\arcs(u)=1$ identically in $x$.
  Integrating $h_b$ over the coordinates outside a separator therefore
  removes their kernel factors exactly.
\item \emph{One kernel everywhere.} The separator marginal of $h_b$ is
  $L_b$ applied to a polynomial in the separator variables of degree at
  most $2r$. Adjacent functionals agree on such polynomials, so adjacent
  separator marginals are \emph{equal as functions}, not merely close.
\end{enumerate}
Running intersection then glues the local laws into one law on the box
(\cref{lem:gluing}). The kernel acts diagonally on Chebyshev polynomials,
multiplying $T_k$ by $\kmult_k$ with $1-\kmult_k\le3k^2/\Dm$, and
\cref{lem:cheb-moment} bounds the Chebyshev moments of $L_b$. Together these
control the change in objective.

\subsection{The Jackson-type kernel}\label{sec:ker-kernel}

Fix an integer $m\ge1$. Let $\tau_j=(m-\abs j)_+$ for $j\in\Z$, and let
\[
  a_k=\sum_{j\in\Z}\tau_j\tau_{j-k},\qquad
  \kmult_k=\frac{a_k}{a_0}\quad(k\ge0),\qquad
  \Dm=2m^2+1 .
\]
The sums are finite, $a_k=a_{-k}$, and $a_k=0$ for $\abs k>2m-2$. Define
the trigonometric polynomial and the bivariate polynomial
\begin{equation}\label{eq:ker-kernel}
  J_m(t)=\frac1{a_0}\Bigl\lvert\sum_{j=0}^{m-1}e^{\mathrm{i}jt}\Bigr\rvert^4,
  \qquad
  \Kj(x,u)=1+2\sum_{k=1}^{2m-2}\kmult_k\,T_k(x)\,T_k(u).
\end{equation}

\begin{lemma}[Jackson-type kernel]\label{lem:jackson}
The following hold.
\begin{enumerate}[label=(\alph*)]
\item $J_m(t)=1+2\sum_{k=1}^{2m-2}\kmult_k\cos(kt)$, and $J_m\ge0$.
\item $a_0=(2m^3+m)/3$, the coefficients $a_k$ are nonnegative integers,
  and $0\le\kmult_k\le1$ for all $k\ge0$, with $\kmult_0=1$ and
  $\kmult_k=0$ for $k>2m-2$.
\item For every integer $k\ge0$,
  \begin{equation}\label{eq:ker-damping}
    0\le1-\kmult_k\le k^2(1-\kmult_1)=\frac{3k^2}{\Dm}.
  \end{equation}
\item $\Kj(\cos\theta,\cos\phi)=\tfrac12\bigl(J_m(\theta-\phi)+J_m(\theta+\phi)\bigr)$.
  Hence $\Kj\ge0$ on $[-1,1]^2$.
\item For every integer $k\ge0$ and every $x\in\R$,
  \begin{equation}\label{eq:ker-diagonal}
    \int_{[-1,1]}T_k(u)\,\Kj(x,u)\,d\arcs(u)=\kmult_k\,T_k(x).
  \end{equation}
  In particular $\int\Kj(x,u)\,d\arcs(u)=1$ identically in $x$, and the
  same holds with the roles of $x$ and $u$ exchanged.
\item $\Kj$ has rational coefficients and degree $2m-2$ in each variable.
\end{enumerate}
\end{lemma}

\begin{proof}
(a) Counting the pairs $(j,j')$ in $\{0,\dots,m-1\}^2$ with $j-j'=\ell$
gives $\abs{\sum_{j=0}^{m-1}e^{\mathrm{i}jt}}^2=\sum_{\ell}\tau_\ell e^{\mathrm{i}\ell t}$.
Squaring, the coefficient of $e^{\mathrm{i}kt}$ is
$\sum_j\tau_j\tau_{k-j}=a_k$, because $\tau$ is even. Dividing by $a_0$ and
using $a_k=a_{-k}$ gives the cosine series. Nonnegativity is clear from the
definition.

(b) Directly, $a_0=\sum_j\tau_j^2=m^2+2\sum_{j=1}^{m-1}j^2=(2m^3+m)/3$. The
$a_k$ are sums of products of nonnegative integers, and $a_k\le a_0$ by the
Cauchy--Schwarz inequality. The support statement was noted above.

(c) For $k\ge1$, the cosine series in (a) gives
$\kmult_k=\int_{-\pi}^{\pi}\cos(kt)J_m(t)\,dt/(2\pi)$; this also holds for
$k>2m-2$, where both sides vanish. Since $J_m\ge0$ has integral one,
$1-\kmult_k=\int(1-\cos kt)J_m(t)\,dt/(2\pi)$. The inequality
$\abs{\sin(k\theta)}\le k\abs{\sin\theta}$ follows by induction from
$\sin((k+1)\theta)=\sin(k\theta)\cos\theta+\cos(k\theta)\sin\theta$. With
$\theta=t/2$ it gives $1-\cos(kt)=2\sin^2(kt/2)\le k^2(1-\cos t)$, and
integration proves the upper bound in \eqref{eq:ker-damping}. For the
equality, $a_0-a_1=\tfrac12\sum_j(\tau_j-\tau_{j-1})^2$. The differences
$\tau_j-\tau_{j-1}$ equal $+1$ for $-m+1\le j\le0$, equal $-1$ for
$1\le j\le m$, and vanish otherwise, so $a_0-a_1=m$ and
$1-\kmult_1=m/a_0=3/(2m^2+1)$.

(d) This follows from (a) and
$\cos(k(\theta-\phi))+\cos(k(\theta+\phi))=2\cos(k\theta)\cos(k\phi)$.

(e) Orthogonality of Chebyshev polynomials under $\arcs$ gives
$\int T_kT_j\,d\arcs=\delta_{jk}/2$ for $k\ge1$ and $\int T_0^2\,d\arcs=1$.
Inserting the definition of $\Kj$ proves \eqref{eq:ker-diagonal} for
$k\le2m-2$, and both sides vanish for $k>2m-2$. The kernel is symmetric.

(f) The $\kmult_k$ are ratios of integers.
\end{proof}

This kernel is the classical squared Fejér, or Jackson-type, construction
used for dense box hierarchies~\cite{dklerk2017-improved-upper-bounds,laurent2023-effective-schmudgen}
and, in tensor form on bags, by Korda, Magron and
Ríos-Zertuche~\cite{korda2025-convergence-rates-sparsity}. We derived the
exact coefficients and estimates used below; we do not claim the kernel as
new. The bound \eqref{eq:ker-damping} holds for every $k$, including
frequencies above the kernel degree, so the objective degree need not be
smaller than the kernel degree.

\subsection{Positive local densities with identical separator marginals}
\label{sec:ker-densities}

For a bag $B$ and $u\in[-1,1]^B$ write
$K_{m,B}(x,u)=\prod_{i\in B}\Kj(x_i,u_i)$. As a polynomial in $x$ it has
total degree at most $2\abs B(m-1)$. Its coefficients are polynomials in
$u$, so $L(K_{m,B}(\cdot,u))$, obtained by applying $L$ in the variable $x$,
is a polynomial in $u$.

\begin{lemma}[Local densities]\label{lem:pre-density}
Let $m\ge1$ and let $(L_b)$ be a feasible preordering family of order $r$
with $v_b(m-1)\le r$ for every bag. For $u\in[-1,1]^{B_b}$ define
\begin{equation}\label{eq:ker-density}
  h_b(u)=L_b\bigl(K_{m,B_b}(\cdot,u)\bigr),\qquad
  h^{(b)}_S(u_S)=L_b\bigl(K_{m,S}(\cdot,u_S)\bigr)\quad(S\subseteq B_b),
\end{equation}
with $h^{(b)}_\varnothing=1$. Then:
\begin{enumerate}[label=(\alph*)]
\item $h_b$ is a polynomial of degree at most $2m-2$ in each coordinate,
  and $h_b(u)\ge0$ for every $u\in[-1,1]^{B_b}$;
\item for every $S\subseteq B_b$,
  $\int h_b\,d\arcs^{B_b\setminus S}=h^{(b)}_S$ identically; in particular
  $\int h_b\,d\arcs^{B_b}=1$;
\item for every edge $e=bc$, $h^{(b)}_{S_e}=h^{(c)}_{S_e}$;
\item for every $\alpha\in\N^{B_b}$ with $\abs\alpha\le2r$,
  \begin{equation}\label{eq:ker-moment-smoothing}
    \int T_\alpha\,h_b\,d\arcs^{B_b}
    =\Bigl(\prod_{i\in B_b}\kmult_{\alpha_i}\Bigr)L_b(T_\alpha).
  \end{equation}
\end{enumerate}
\end{lemma}

\begin{proof}
(a) The degree statement is clear. Fix $u$. By \cref{lem:jackson}(d) the
polynomial $\Kj(\cdot,u_i)$ is nonnegative on $[-1,1]$ and has degree at
most $2(m-1)$. \Cref{lem:interval-sos} with $k=m-1$ gives
$\Kj(x_i,u_i)=\sigma_{0,i}(x_i)+(1-x_i^2)\sigma_{1,i}(x_i)$ with
$\sigma_{0,i}\in\sos_{2m-2}$ and $\sigma_{1,i}\in\sos_{2m-4}$ (for $m=1$,
$\Kj=1$ and $\sigma_{1,i}=0$). Multiplying over $i\in B_b$,
\[
  K_{m,B_b}(x,u)=\sum_{I\subseteq B_b}
  \Bigl(\prod_{i\notin I}\sigma_{0,i}(x_i)\prod_{i\in I}\sigma_{1,i}(x_i)\Bigr)
  \prod_{i\in I}(1-x_i^2).
\]
The bracket is a sum of squares of polynomials of degree at most
$v_b(m-1)-\abs I\le r-\abs I$, so every term belongs to
$\Pre_r(B_b)$, and $h_b(u)\ge0$. The representation may depend on $u$ in an
arbitrary way; it is used only at the fixed point $u$.

(b) Every coefficient of $K_{m,B_b}(x,u)$ as a polynomial in $x$ is a
polynomial in $u$, and $L_b$ is linear, so integration in $u$ and
application of $L_b$ commute. By \cref{lem:jackson}(e), integrating
$\Kj(x_i,u_i)$ in $u_i$ gives the constant one. Integrating out the
coordinates of $B_b\setminus S$ therefore leaves $L_b(K_{m,S}(\cdot,u_S))$.

(c) $K_{m,S_e}(x,u)$ is a polynomial in $x_{S_e}$ of total degree at most
$2\abs{S_e}(m-1)\le2r$. Exact separator consistency
(\cref{def:sparse-hierarchy}(iii)) applied coefficientwise in $u$ gives the
claim.

(d) Expanding the definition, $K_{m,B_b}(x,u)=\sum_\beta
\bigl(\prod_ic_{\beta_i}\kmult_{\beta_i}\bigr)T_\beta(x)T_\beta(u)$, where
$\beta$ ranges over multi-indices with $\beta_i\le2m-2$, $c_0=1$ and
$c_k=2$ for $k\ge1$. Since $\int T_{\alpha_i}T_{\beta_i}\,d\arcs=
\delta_{\alpha_i\beta_i}/c_{\alpha_i}$, integration against $T_\alpha(u)$
gives $(\prod_i\kmult_{\alpha_i})L_b(T_\alpha)$ if $\alpha_i\le2m-2$ for
all $i$, and zero otherwise; in the latter case some $\kmult_{\alpha_i}$
vanishes. The value $L_b(T_\alpha)$ is defined because $\abs\alpha\le2r$.
\end{proof}

\begin{remark}[Degree accounting and the role of the preordering]
\label{rem:ker-degree}
Two features of \cref{lem:pre-density} cannot be weakened casually.

First, positivity is certified through the univariate certificates, not
through the total degree of the tensor kernel. A product of interval
nonnegative polynomials whose total degree is at most $2r$ can be negative
under a feasible functional. Take two variables, $r=1$, and
$L(1)=1$, $L(x)=L(y)=\tfrac12$, $L(x^2)=L(y^2)=1$, $L(xy)=-\tfrac12$. The
moment matrix on $(1,x,y)$ has eigenvalues $0,\tfrac32,\tfrac32$, and
$L(1-x^2)=L(1-y^2)=0$, so $L$ is $\Pre_1$-positive (at order one the
preordering has no products of generators). Yet
$L((1-x)(1-y))=-\tfrac12$, although $(1-x)(1-y)\ge0$ on the box and has
degree two. A univariate factor of degree $q$ needs interval certificate
degree $2\lceil q/2\rceil$; the even kernel degree $2(m-1)$ makes the count
in \cref{lem:pre-density}(a) exact.

Second, the expansion of $K_{m,B_b}$ contains products of distinct box
generators. The ordinary module does not contain these products, and a
normalized kernel that is merely nonnegative on the interval does not in
general have an ordinary-module representation of the tensor product. This
is why \cref{sec:ordinary} uses a different kernel.

Finally, the same kernel must be used for a given coordinate in all bags
containing it, with the same parameter $m$. With different kernels in
two bags, equal input moments would not give equal output marginals.
\end{remark}

\subsection{The rounding theorem}\label{sec:ker-theorem}

Recall that $d$ is the largest total degree of the local polynomials $f_b$
and that $\Abud=\sum_b\Abud(f_b)$ is the weighted Chebyshev budget of
\eqref{eq:set-budgets}; it is zero for constant local terms and is not
affected by additive constants.

\begin{theorem}[Rounding for the sparse box preordering]\label{thm:pre}
Let $r\ge\max\{w,d\}$ and $m=\lfloor r/w\rfloor+1$. For every feasible
preordering family $(L_b)$ of order $r$ there is a Borel probability
measure $\nu$ on $[-1,1]^n$ such that the marginal of $\nu$ on
$[-1,1]^{B_b}$ is $h_b\,d\arcs^{B_b}$ for every bag, with $h_b$ from
\eqref{eq:ker-density}, and
\begin{equation}\label{eq:ker-rounding}
  \Bigl\lvert\int f\,d\nu-\sum_bL_b(f_b)\Bigr\rvert\le\frac{3\Abud}{\Dm}.
\end{equation}
Consequently
\begin{equation}\label{eq:ker-gap}
  0\le\fstar-\rhopre_r\le\frac{3\Abud}{2m^2+1}\le\frac{3w^2\Abud}{2r^2},
\end{equation}
and for every feasible family there is a point $x\in[-1,1]^n$ with
$f(x)\le\sum_bL_b(f_b)+3\Abud/\Dm$.
\end{theorem}

\begin{proof}
Since $m-1=\lfloor r/w\rfloor$, we have $v_b(m-1)\le w\lfloor r/w\rfloor\le r$
for every bag, so \cref{lem:pre-density} applies. By parts (a) and (b), each
$\nu_b=h_b\,d\arcs^{B_b}$ is a probability measure on $[-1,1]^{B_b}$. By
parts (b) and (c), adjacent laws have the same separator marginal, namely
$h^{(b)}_{S_e}\,d\arcs^{S_e}$. \Cref{lem:gluing} gives a probability measure
$\nu$ on $[-1,1]^n$ with marginals $\nu_b$, and
$\int f\,d\nu=\sum_b\int f_b\,d\nu_b$.

Fix a bag and expand $f_b=\sum_\alpha c_{b,\alpha}T_\alpha$; every
$\alpha$ with $c_{b,\alpha}\ne0$ has $\abs\alpha\le d\le r$. By
\eqref{eq:ker-moment-smoothing},
\[
  \int f_b\,d\nu_b-L_b(f_b)=\sum_\alpha c_{b,\alpha}
  \Bigl(\prod_i\kmult_{\alpha_i}-1\Bigr)L_b(T_\alpha).
\]
\Cref{lem:cheb-moment} gives $\abs{L_b(T_\alpha)}\le1$. For numbers
$\kmult_i\in[0,1]$ we have $0\le1-\prod_i\kmult_i\le\sum_i(1-\kmult_i)$ by
telescoping, so \eqref{eq:ker-damping} gives
$\abs{1-\prod_i\kmult_{\alpha_i}}\le3\sum_i\alpha_i^2/\Dm$. Hence
$\abs{\int f_b\,d\nu_b-L_b(f_b)}\le3\Abud(f_b)/\Dm$. Summing over bags proves
\eqref{eq:ker-rounding}.

Since $\nu$ is supported on the box, $\fstar\le\int f\,d\nu\le
\sum_bL_b(f_b)+3\Abud/\Dm$, and some point of the box has value at most
$\int f\,d\nu$. Taking the infimum over feasible families gives the first
inequality in \eqref{eq:ker-gap} together with
\cref{lem:weak-duality}. Finally $m>r/w$, so $2m^2+1>2r^2/w^2$.
\end{proof}

The theorem is a statement about \emph{every} feasible family, not only an
optimal one. By \cref{lem:compact} an optimal family exists; rounding it
gives a law $\nu$ with $\int f\,d\nu\le\rhopre_r+3\Abud/\Dm$. The objective
value of an arbitrary feasible family is not a lower bound on $\fstar$; the
moment value $\rhopre_r$ is. The constant depends on the chosen
decomposition of $f$ through $\Abud$, and $\Abud$ can grow with the number
of bags. The rounded law is supported on the box, but it need not satisfy
any further constraint, and it need not respect prescribed integer values.

\begin{remark}[Explicit gluing for polynomial densities]\label{rem:ker-gluing}
For the laws in \cref{thm:pre}, the conditional laws used in
\cref{lem:gluing} can be written down, so no general disintegration theorem
is needed. Root $\Tree$ and let bag $b$ have parent separator $S$ and new
coordinates $C=B_b\setminus S$. Let $h_S=h^{(b)}_S$. Define the kernel
$\kappa(u_S,du_C)=h_b(u_S,u_C)h_S(u_S)^{-1}\,d\arcs^C(u_C)$ where
$h_S(u_S)>0$, and $\kappa(u_S,\cdot)=\arcs^C$ where $h_S(u_S)=0$. For
measurable $A\subseteq[-1,1]^S$ and $F\subseteq[-1,1]^C$, the parent law
gives $A$ the mass $\int_Ah_S\,d\arcs^S$, so the glued law gives $A\times F$
the mass $\int_{A\cap\{h_S>0\}}\int_Fh_b\,d\arcs^C\,d\arcs^S$. Since
$h_b\ge0$ and $\int h_b\,d\arcs^C=h_S$, the inner integral vanishes where
$h_S=0$, so this mass equals $\nu_b(A\times F)$. Rectangles determine
$\nu_b$, and the induction of \cref{lem:gluing} goes through. The same
construction with sums applies to the finite laws in \cref{cor:pre-grid},
and with $h_b^+$ in place of $h_b$ to the laws of \cref{sec:ordinary}.
\end{remark}

\subsection{Finite grid rounding and certificates}\label{sec:ker-grid}

Because each density $h_b$ is a polynomial of known coordinate degree, the
integrals in \cref{thm:pre} can be replaced by an exact quadrature on a
fixed Chebyshev grid. The rounded law then lives on a finite grid, and the
best grid point can be found by dynamic programming over the junction tree.

\begin{lemma}[Chebyshev quadrature]\label{lem:ker-quadrature}
Let $N\ge1$ and $\xi_j=\cos\bigl((2j-1)\pi/(2N)\bigr)$ for
$j=1,\dots,N$. For every univariate polynomial $q$ of degree at most $2N-1$,
$\frac1N\sum_{j=1}^Nq(\xi_j)=\int q\,d\arcs$.
\end{lemma}

\begin{proof}
It suffices to take $q=T_k$, $0\le k\le2N-1$. For $k=0$ both sides equal
one. For $1\le k\le2N-1$, put $\theta=k\pi/(2N)\in(0,\pi)$. Then
$\sum_{j=1}^Ne^{\mathrm{i}(2j-1)\theta}=e^{\mathrm{i}\theta}
(1-e^{2\mathrm{i}N\theta})/(1-e^{2\mathrm{i}\theta})$, and the denominator
is nonzero. If $k$ is even, $e^{2\mathrm{i}N\theta}=e^{\mathrm{i}k\pi}=1$
and the sum vanishes. If $k$ is odd, the sum equals
$2e^{\mathrm{i}\theta}/(1-e^{2\mathrm{i}\theta})=\mathrm{i}/\sin\theta$,
whose real part is zero. Hence $\sum_jT_k(\xi_j)=0=N\int T_k\,d\arcs$.
\end{proof}

Let $d_\infty$ be the largest coordinate degree of the local polynomials,
that is, the largest $\alpha_i$ with $c_{b,\alpha}\ne0$ (zero if all $f_b$
are constant).

\begin{corollary}[Rounding to a fixed grid]\label{cor:pre-grid}
Let $r$ and $m$ be as in \cref{thm:pre}, let $N_{\mathrm g}=m+\lfloor
d_\infty/2\rfloor$, and let $\Xi=\{\xi_1,\dots,\xi_{N_{\mathrm g}}\}$ be the
nodes of \cref{lem:ker-quadrature} with $N=N_{\mathrm g}$. For every
feasible preordering family $(L_b)$ of order $r$,
\begin{equation}\label{eq:ker-grid}
  \min_{x\in\Xi^n}f(x)\le\sum_bL_b(f_b)+\frac{3\Abud}{\Dm}.
\end{equation}
In particular $\fstar\le\min_{\Xi^n}f\le\rhopre_r+3\Abud/\Dm\le
\fstar+3\Abud/\Dm$. Given the tables of local values $f_b$ on
$\Xi^{B_b}$, a grid minimizer can be computed by dynamic programming over
$\Tree$ with $\sum_bN_{\mathrm g}^{v_b}$ table entries and at most
$2tN_{\mathrm g}^w$ additions and comparisons.
\end{corollary}

\begin{proof}
Since $2\lfloor d_\infty/2\rfloor\ge d_\infty-1$, we have
$2N_{\mathrm g}-1\ge2(m-1)+d_\infty$. The density $h_b$ has degree at most
$2m-2$ in each coordinate, and $f_bh_b$ has degree at most $2m-2+d_\infty$
in each coordinate. Applying \cref{lem:ker-quadrature} coordinate by
coordinate, the tensor rule with weights $N_{\mathrm g}^{-v_b}$ integrates
$h_b$, $f_bh_b$, and $h_b$ with respect to any subset of its coordinates
exactly. Define finite bag laws
$p_b(\xi_{j})=N_{\mathrm g}^{-v_b}h_b(\xi_j)$ for
$\xi_j\in\Xi^{B_b}$. They are nonnegative by \cref{lem:pre-density}(a) and
have total mass one. Summing $p_b$ over the coordinates of $B_b\setminus S$
gives $N_{\mathrm g}^{-\abs S}h^{(b)}_S$ at the remaining nodes, by exact
quadrature and \cref{lem:pre-density}(b); by part (c) adjacent bags have
equal separator marginals. \Cref{lem:gluing} for finite sets (or
\cref{rem:ker-gluing} with sums) gives a law on $\Xi^n$ with these bag
marginals. Its expected objective is
$\sum_b\sum_jp_b(\xi_j)f_b(\xi_j)=\sum_b\int f_bh_b\,d\arcs^{B_b}$, which is
the quantity bounded in the proof of \cref{thm:pre}. Some grid point is no
worse than the mean, proving \eqref{eq:ker-grid}. The chain of inequalities
follows from \cref{lem:weak-duality} and \eqref{eq:ker-gap}, applied to an
optimal family (\cref{lem:compact}).

For the dynamic program, root $\Tree$ and process bags from the leaves up.
For a bag $b$ with parent separator $S$, form the table of $f_b$ plus the
messages received from its children, each a function of a child separator
and hence a table lookup, and minimize over the coordinates of
$B_b\setminus S$ to obtain the message to the parent. This uses
$(1+\#\mathrm{children}(b))N_{\mathrm g}^{v_b}$ additions and comparisons
at bag $b$, and $\sum_b(1+\#\mathrm{children}(b))=2t-1$. Storing the
minimizing choices recovers a grid minimizer.
\end{proof}

The operation count is in exact real arithmetic; the nodes $\xi_j$ are
algebraic numbers, and no bit-complexity claim is made. The grid bound
$\min_{\Xi^n}f-\fstar\le3\Abud/\Dm$ on its own is a polynomial
approximation statement of the type established by Piazzon and
Vianello~\cite{piazzon2018-chebyshev-grids}. The proof describes the dynamic
programming algorithm that minimizes the grid objective on a junction tree. Grid
optimization with dynamic programming therefore attains the same exponent
without solving any SDP. The content of \cref{cor:pre-grid} is the
comparison: the grid optimum is at most the objective of \emph{every}
feasible point of the SDP plus $3\Abud/\Dm$. For an optimal SDP point,
the grid optimum and the moment value differ by at most this error.
Bounded-treewidth polynomial
optimization also admits LP approximation schemes~\cite{bienstock2018-lp-formulations-for-polynomial-optimization},
so no algorithmic advance is claimed.

\begin{corollary}[Sparse preordering certificate]\label{cor:pre-certificate}
Let $r$ and $m$ be as in \cref{thm:pre}. Then
\[
  f-\fstar+\frac{3\Abud}{\Dm}\in\sum_{b=1}^t\Pre_r(B_b),
\]
with real coefficients and every summand of total degree at most $2r$.
Equivalently, $\lampre_r=\rhopre_r\ge\fstar-3\Abud/\Dm$.
\end{corollary}

\begin{proof}
By \cref{lem:slater}, $f-\rhopre_r=\sum_bq_b$ with
$q_b\in\Pre_r(B_b)$. By \cref{thm:pre}, the constant
$\rhopre_r-\fstar+3\Abud/\Dm$ is nonnegative; add it to $q_1$, where it is a
square of a constant.
\end{proof}

The certificate has real coefficients. Rational certificates and their size
are treated in \cref{sec:certificates}. Finite SDP dual attainment, used
here, is different from the existence of exact polynomial separator
functions; \cref{thm:quad-sharp} shows that the latter can fail for a fixed
quadratic objective.

\subsection{Relation to prior work}\label{sec:ker-prior}

\begin{remark}[What is prior in \cref{thm:pre}]\label{rem:pre-prior}
The rate is not new. Magron's lecture slides of 7 July
2025~\cite{magron2025slides-lorentz} (slide 23 of 44) and of 16 February
2026~\cite{magron2026slides-tenors} (slide 35 of 90) state an $O(d^{-2})$
gap for two overlapping bags with local full preorderings and total-degree
truncation $2d$, attributed to Korda, Ríos-Zertuche and Magron. The
published article of these
authors~\cite{korda2025-convergence-rates-sparsity}, Theorem~6, gives for
the sparse box preordering a degree bound equivalent to the rate
$O(r^{-2/(w+3)})$ for fixed data, using coordinatewise degree bounds; the
difference from total degree costs only width-dependent factors. We could
not resolve the difference between the slides and the published theorem,
and we do not assume either is in error. The published proof first
approximates separator value functions by polynomials to obtain positive
bag polynomials and then applies a sparse Jackson operator; the tensor
Jackson kernel and its preordering positivity already appear there.

Dense box preordering hierarchies have error $O(r^{-2})$ by Laurent and
Slot~\cite{laurent2023-effective-schmudgen}, using the same kind of kernel,
which de Klerk, Hess and Laurent~\cite{dklerk2017-improved-upper-bounds}
used earlier for upper-bound hierarchies. Dense full preorderings on
products of standard domains are treated by
Magron~\cite{magron2026-convergence-rates-for-polynomial-optimization}. For
actual measures on the torus, equal low-order Fourier moments give equal
Jackson convolutions, as observed by Catala, Hockmann, Kunis and
Wageringel~\cite{catala2024-singular-measures}. Applying the cosine map gives
the corresponding statement for Chebyshev moments. Gluing consistent local
measures under running intersection is classical in sparse
optimization~\cite{lasserre2006-convergent-sdprelaxations-in-polynomial-optimization}.
Gamertsfelder and Mourrain~\cite{gamertsfelder2025-countable-gmp} transfer
effective positivity rates to generalized moment relaxations under an
assumption that a dual with finitely supported, here polynomial, separator
multipliers is attained; \cref{sec:sharpness} gives a quadratic example in
which no exact polynomial separator decomposition exists.

What \cref{thm:pre} adds, to our knowledge, is a direct finite-order
primal construction: every feasible sparse preordering family is turned
into exactly consistent local probability densities, with the explicit
constant $3\Abud/\Dm$ and without approximating separator value functions
or assuming dual attainment. Independently approximating each local
functional by a measure, as in fixed-degree truncated moment
approximation~\cite{tran2026-truncated-moment-sequences}, does not by itself
produce equal separator marginals.
\end{remark}
